%% Native common Q-Res calculus with complete finite-PCNF proofs.
\chapter{Q-Resolution and Quantified Certificates}
\label{sec:p12-q-resolution-and-quantified-certificates}

Definition~\ref{def:p04-qbf-languages} and
Theorem~\ref{thm:p04-tqbf-pspace} already fix the QBF language and prove its
PSPACE completeness.  This chapter supplies a different object: a quantified
clause refutation and its strategy-relative checking guarantee.  It uses the
actual common Q-Res presentation in the July~2021 author manuscript
\cite[Sec.~2.2, Fig.~2.1]{BeyersdorffEtAl2021QBFJulyAuthor}, keeping its
separate identity from the final handbook chapter and the original 1995
paper to which it attributes the calculus.  The independent complete proofs
below concern precisely the displayed finite-PCNF rules.  They do not identify
all later resolution, dependency, expansion or QRAT systems with Q-Res.

\section{Prefix order and legal strategy dependencies}

\begin{hbdefinition}{def:p12-qres-pcnf}{The native PCNF interface}
A closed prenex conjunctive-normal-form input is \(\phi=\Pi\psi\), with
finite prefix \(\Pi=Q_1X_1\cdots Q_kX_k\),
\(Q_i\in\{\exists,\forall\}\), and disjoint finite variable blocks
\(X_i\).  Adjacent nonempty blocks have different quantifiers, and every
matrix variable is in the prefix.  An empty prefix is permitted for the
quantifier-free terminal case; unused variables have no effect.  The matrix
\(\psi\) is a finite set of clauses, each a finite set of signed literals.
Complement \(\bar l\) flips a literal's sign, with
\(\overline{\bar l}=l\), rather than adding a double-negation formula.
A clause containing both signs of one variable is tautological.  The empty
clause \(\varnothing\), also denoted \(\bot\) in the source, is false;
the empty matrix is true.

A Boolean valuation \(v\) on a specified finite variable set maps
exactly that set to \(\{0,1\}\). In the elimination proof its domain
is the remaining earlier variables: ``complete'' refers to that stated
finite domain, not to every variable of the surrounding propositional
language. A complete play instead assigns all variables of the input prefix.

For \(x\in X_i\), write \(\operatorname{lv}(x)=i\) and
\(\operatorname{quant}(\Pi,x)=Q_i\), extended to literals by their
underlying variable.  The source order \(x\leq_\Pi y\) compares block
levels.  In particular an existential literal \(e\leq_\Pi u\) with
universal \(u\) is in a strictly earlier block, since differently quantified
variables cannot share a block.  This is the linear prefix restriction, not
an arbitrary declared dependency relation
\cite[Sec.~2.2]{BeyersdorffEtAl2021QBFJulyAuthor}.
The source uses \(\operatorname{var}(l)\) for a literal's variable and
\(\operatorname{var}(\psi)\) for all matrix variables.  Its rule figure
writes \(\operatorname{vars}(\Pi)\) for the prefix's variable set, the
union of its blocks.  The sheet retains that locally typed figure notation.
\end{hbdefinition}

\begin{hbdefinition}{def:p12-qres-strategies}{Skolem and Herbrand strategy certificates}
The quantifiers are interpreted recursively: \(\exists x\) accepts one of
the two values, while \(\forall x\) requires both.  Equivalently the players
choose values in prefix order, and the existential player wins exactly when
the final matrix is true.  A legal existential strategy assigns each
existential variable \(e\) by a Boolean function \(f_e\) of only the
universal variables preceding it.  A legal universal strategy assigns each
universal \(u\) by a Boolean function \(g_u\) of only the preceding
existential variables.  The players' own earlier values are recovered from
their earlier functions and need not be extra independent arguments.
These are respectively Skolem and Herbrand functions in the source's
strategy presentation \cite[Sec.~2.2]{BeyersdorffEtAl2021QBFJulyAuthor}.

A strategy certificate specifies each allowed argument list and the complete
Boolean table for its function.  An empty argument list has one constant
value.  A winning existential certificate must make \(\psi\) true for
every universal assignment; a winning universal certificate must make it
false for every existential assignment.  In an alternating prefix a single
complete valuation records only one play and is not in general a strategy
certificate.  Constant strategies are permitted, but must still pass every
opponent assignment; a purely existential input reduces to ordinary valuation
checking.  An existential function of a later universal variable is illegal
even if its values make the matrix true.
\end{hbdefinition}

\begin{hbtheorem}{thm:p12-qres-game-strategies}{proved=annex:p12-qres-game-strategies}{Finite game semantics and strategy construction}
For every closed finite PCNF, exactly one player has a winning legal strategy.
The existential player has one precisely when the QBF is true; the universal
player has one precisely when it is false.  Backward evaluation of the finite
assignment tree constructs its finite strategy tables with only the permitted
preceding-opponent dependencies, including empty prefixes and blocks.
\end{hbtheorem}

\section{The exact common Q-Res rule sheet}

% Exact common Q-Res rule presentation: July 2021 author manuscript, Fig. 2.1.
% This is not a claim of source-byte reconstruction of the unheld 1995 paper.
\calculussheet{sheet:p12-q-resolution}{Q-Resolution: the common Q-Res presentation}{Language: closed PCNF $\Pi\psi$ with finite prefix blocks; proof lines: non-tautological clauses over prefix variables; negation: literal complement $\bar l$; equality: none. Prefix and matrix remain fixed during checking.}
\hbrule{rule:p12-q-resolution-init}{cl-init}{}{C}{$C\in\psi$ and $\{x,\bar x\}\not\subseteq C$ for every $x\in\operatorname{vars}(\Pi)$}
\hbrule{rule:p12-q-resolution-red}{red}{C\cup\{l\}}{C}{$\operatorname{quant}(\Pi,l)=\forall$; the parent is non-tautological; $l'\leq_\Pi l$ for every $l'\in C$ with $\operatorname{quant}(\Pi,l')=\exists$}
\hbrule{rule:p12-q-resolution-res}{res}{C_1\cup\{p\}\qquad C_2\cup\{\bar p\}}{C_1\cup C_2}{$\operatorname{quant}(\Pi,p)=\exists$; $\bar p\notin C_1$, $p\notin C_2$; $\{x,\bar x\}\not\subseteq C_1\cup C_2$ for every $x\in\operatorname{vars}(\Pi)$}


All proof lines are interpreted with respect to the same fixed \(\Pi\).
The axiom \texttt{cl-init} selects a non-tautological input clause.
\texttt{red} removes a trailing universal literal only when every existential
literal of the conclusion is earlier in the prefix.  If the conclusion has
no existential literal, this order condition is vacuous, but the parent must
still be non-tautological.  \texttt{res} has an existential pivot and rejects
a tautological conclusion.  Removing a blocked universal or resolving on a
universal pivot is not an instance of this sheet.  The latter restriction
identifies this calculus; it is not a claim that every separately specified
universal-pivot calculus is unsound.

\begin{hbdefinition}{def:p12-qres-strategy-consequence}{What a derived clause preserves}
A legal existential strategy \(S\) satisfies a clause \(C\) when every
universal assignment, completed with the functions of \(S\), makes \(C\)
true.  A proof invariant says that every winning strategy for the input
satisfies each derived clause.  It does \emph{not} say that each derived
clause is a propositional consequence of the matrix alone.  For example,
\(a\) can be obtained from \(a\lor u\) under prefix \(\exists a\forall u\)
by universal reduction, although the valuation \(a=0,u=1\) satisfies the
parent and falsifies the conclusion.  The prefix is part of the inference.
\end{hbdefinition}

\begin{hbtheorem}{thm:p12-qres-rule-soundness}{proved=annex:p12-qres-rule-soundness}{Rule and refutation soundness with the full prefix}
Each displayed rule preserves the invariant of
Definition~\ref{def:p12-qres-strategy-consequence}.  Consequently deriving
the empty clause proves the original QBF false.  Universal reduction
preserves a winning strategy because changing its removed universal value
cannot change any remaining existential literal, whose dependencies are
earlier.  The proof includes the non-tautology condition, the case with no
remaining existential literal, and vacuous prefix variables.
\end{hbtheorem}

\begin{hbexample}{def:p12-qres-native-derivation}{The source's blocked-then-legal reduction}
Retain the source PCNF
\[
 \exists x_1,x_3,x_4\;\forall y_5\;\exists x_2\;
 \{\{\bar x_1,x_2\},\{x_3,y_5,\bar x_2\},
   \{x_4,\bar y_5,\bar x_2\},\{\bar x_3,\bar x_4\}\}.
\]
Resolving the last and third clauses on \(x_4\) gives
\(C_0=\{\bar x_3,\bar y_5,\bar x_2\}\).  Reduction of \(\bar y_5\)
is forbidden because \(x_2\) follows \(y_5\).  Resolving \(C_0\) with
\(\{\bar x_1,x_2\}\) on \(x_2\) gives
\(C_1=\{\bar x_1,\bar x_3,\bar y_5\}\).  Now both existential variables
are earlier, so reduction yields \(C_2=\{\bar x_1,\bar x_3\}\)
\cite[Example~2.3.1]{BeyersdorffEtAl2021QBFJulyAuthor}.
This is a derivation, not a refutation: the input is true, witnessed by
constant-zero existential functions.  The derived clause constrains winning
strategies and does not justify changing the source example into an UNSAT
claim.
\end{hbexample}

\section{A general certificate construction, not a SAT replacement}

\begin{hbmethod}{def:p12-qres-elimination}{Right-to-left quantified elimination}
Choose an order within each block and expand the finite prefix into scalar
quantifiers.  Work from its innermost remaining variable outward, maintaining
a CNF over only the still-uneliminated variables.  Remove tautological input
clauses, which are semantically true, and record the others as axioms.
For innermost universal \(u\), remove its literal from every containing
clause; the side condition of \texttt{red} holds since no later existential
variable remains in those clauses.  Keep clauses not containing \(u\).

For innermost existential \(x\), partition clauses into those not mentioning
\(x\), those \(A\cup\{x\}\), and those \(B\cup\{\bar x\}\).
Replace the last two groups by every non-tautological \(A\cup B\) from
a positive/negative pair, recording its resolution parents.  Ignore a
tautological resolvent as true, rather than inserting it into the proof.
The working matrix drops the parent clauses involving the eliminated variable;
the proof record keeps their earlier identifiers.  This is an elimination
construction with a new working matrix, not a proof-log deletion rule.
No new variable or forbidden strategy dependency is introduced.
\end{hbmethod}

\begin{hblemma}{thm:p12-qres-elimination-step}{proved=annex:p12-qres-elimination-step}{Exact innermost existential and universal elimination}
For every assignment to the remaining earlier variables, the existential
replacement computes exactly \(\exists x\) of the old matrix, and the
universal replacement computes exactly \(\forall u\) of it.  The
tautological-resolvent omission and non-tautological-parent restriction have
their actual semantic meanings in these equalities.  Hence applying the
unchanged earlier prefix preserves QBF truth at each step.
\end{hblemma}

\begin{hbtheorem}{thm:p12-qres-completeness}{proved=annex:p12-qres-completeness}{Constructive refutational completeness}
For every false closed finite PCNF, the elimination construction terminates
and produces a valid Q-Res derivation of the empty clause under the original
prefix.  Conversely a valid such derivation exists only for a false input.
This is general finite-PCNF completeness, with no bound on alternations,
clause width or variable count.  It makes no polynomial certificate-size
or practical solver-termination claim.
\end{hbtheorem}

\begin{hbexample}{def:p12-qres-alternating-certificate}{An alternating refutation and a legal counterstrategy}
Take prefix \(\exists a\forall u\exists b\) and four input clauses
\(\{a,u,b\}\), \(\{a,u,\bar b\}\),
\(\{\bar a,\bar u,b\}\), \(\{\bar a,\bar u,\bar b\}\).
The complete certificate has the following clause records; the parent numbers
refer only to earlier records:
\begin{center}
\begin{tabular}{@{}lll@{}}
\toprule
Record & Clause & Justification \\
\midrule
1--4 & The four displayed clauses & \texttt{cl-init} \\
5 & \(\{a,u\}\) & \texttt{res} on \(b\), parents 1,2 \\
6 & \(\{\bar a,\bar u\}\) & \texttt{res} on \(b\), parents 3,4 \\
7 & \(\{a\}\) & \texttt{red} of \(u\), parent 5 \\
8 & \(\{\bar a\}\) & \texttt{red} of \(\bar u\), parent 6 \\
9 & \(\varnothing\) & \texttt{res} on \(a\), parents 7,8 \\
\bottomrule
\end{tabular}
\end{center}
Initially \(u\) is blocked by later \(b\); its reduction becomes legal
only after the two resolutions.  The universal counterstrategy is
\(g_u(a)=a\), a function of an earlier existential value.  If \(a=u=0\)
the first pair requires both values of \(b\); if \(a=u=1\) the second pair
does.  No later choice of \(b\) can win.  The propositional matrix itself
has models, for example \(a=0,u=1,b=0\), so ordinary matrix SAT is not
the conclusion certified here.
\end{hbexample}

\section{Checking and strategy certificates}

\begin{hbmethod}{def:p12-qres-proof-checker}{The abstract quantified proof checker}
Retain the exact input prefix/matrix.  Each record gives its clause, rule tag,
earlier parent identifiers, and selected pivot or removed literal where
required.  Check closure and quantifier types; check every literal-set equality,
input membership and non-tautology guard; and check every prefix-order test.
Accept refutation only if a checked record is the empty clause.  Forward
references, a changed prefix, a failed rule condition or an unchecked final
empty record reject the proof.  All these are finite syntactic checks and
are polynomial in the explicit input plus certificate length.  This
description is a checker specification, not an executed or verified binary.
The logical soundness theorem supplies its guarantee; parser and encoding
correspondence remain separately required.
\end{hbmethod}

\begin{hbmethod}{def:p12-qres-strategy-checker}{Finite direct strategy checking}
For each controlled variable check that its table is total on exactly its
declared preceding-opponent argument list.  Reject later-variable arguments.
Enumerate every assignment to the opposing player's variables.  Evaluate the
controlled functions and the matrix on the induced complete assignment.
Require all matrix values true for an existential certificate, or all false
for a universal certificate.  An incomplete table or one good play is not
acceptance.  Enumeration can be exponential in the variable count; this is
not a claim that succinct arbitrary circuit strategies are polynomially
checkable merely because their circuit is small.
\end{hbmethod}

\begin{hbtheorem}{thm:p12-qres-strategy-checking}{proved=annex:p12-qres-strategy-checking}{Exact meaning of finite strategy acceptance}
The direct finite checker terminates and accepts precisely the legal winning
table strategies.  Together with the finite-game construction it supplies
strategy certificates for true and false inputs.  A Q-Res refutation is a
false-input certificate; a true-input strategy is not an empty-clause proof.
No polynomial-time strategy-extraction algorithm from an arbitrary Q-Res log
is asserted by this enumeration construction.
\end{hbtheorem}

\section{Failed rule conditions and separate presentations}

\begin{hbexample}{def:p12-qres-illegal-reduction}{Illegal reduction and tautological resolution}
The source's true formula is
\[
 \forall x\exists y\;\{\{x,\bar y\},\{\bar x,y\}\}.
\]
Its winning strategy is \(f_y(x)=x\).  Removing \(x\) from the first
clause or \(\bar x\) from the second is illegal: the later existential
\(y\) is allowed to respond to \(x\).  Those wrong reductions would give
\(\{\bar y\}\) and \(\{y\}\) and a false refutation.  Likewise resolving
the two inputs on \(y\) gives \(\{x,\bar x\}\), a tautology forbidden
by \texttt{res}; reducing both universal literals from it also violates
\texttt{red}'s non-tautology premise
\cite[Example~2.3.2]{BeyersdorffEtAl2021QBFJulyAuthor}.
Neither an arbitrary independence declaration nor treating a tautology as an
ordinary reduction parent repairs the certificate.
\end{hbexample}

Reversing the prefix gives
\(\exists y\forall x\;\{\{x,\bar y\},\{\bar x,y\}\}\), which is
false: \(y\) must be constant, and the universal player chooses
\(x=1-y\).  The apparent existential strategy \(y=x\) now reads a later
universal variable and is forbidden.  Thus equal matrices do not determine
equal quantified problems or equal strategy certificates.

\begin{hbdefinition}{def:p12-qres-full-reduction}{The distinct fully reduced convention}
The 2013 ECCC presentation assumes input clauses already fully universally
reduced and defines a Q-resolvent by first forming a non-tautological
propositional resolvent and then removing all unblocked universal literals
\cite[Preliminaries, p.~2]{JanotaMarquesSilva2013QBFExpansionECCC97}.
Write \(\operatorname{ur}_{\Pi}(C)\) for this full reduction.  This convention
packages several source-sheet \texttt{red} steps into one operation; it is
not a claim that the later presentation is byte-identical to the
historical 1995 publication. The input normalization itself has a quantified
truth-preservation obligation, not merely a renaming of clauses.
\end{hbdefinition}

\begin{hbtheorem}{thm:p12-qres-normalized-bridge}{proved=annex:p12-qres-normalized-bridge}{The full-reduction convention at its actual standing}
Full reduction is a finite sequence of legal \texttt{red} steps and preserves
exactly the winning existential strategies of the clause matrix.  Every
fully reduced resolution step expands into source-sheet resolution followed
by legal reduction steps.  Both conventions admit a refutation exactly for
false closed finite PCNFs, after the declared input normalization.  This is
a semantic/refutational bridge, not linewise equality of every intermediate
clause, whole-source equivalence or a proof-size comparison.
\end{hbtheorem}

Long-distance Q-resolution uses specially governed merged universal literals;
QU-resolution admits universal pivots; dependency-aware reduction replaces
linear order tests with separately justified dependency information.  They
have distinct syntax, side conditions and proof obligations
\cite[Sec.~2.3.2]{BeyersdorffEtAl2021QBFJulyAuthor}.  The ordinary tautology
in the failed example is not such a merged-literal object.  Clause/cube
QCDCL, expansion proofs, QRAT and strategy-extraction constructions also
remain separate from the common Q-Res sheet.  None is proved by importing
ordinary SAT resolution, nor are QBF prefix dependencies preserved merely
because a preprocessing step preserves propositional matrix satisfiability.
